% ==================================================
% Chapter EVALUATION
% ==================================================

\section{测评细则}
\subsection{电流计部分}
\begin{enumerate}
    \item 优先测试电流转换功能，使用台式供电，共模电压在要求范围内，输出对应的电流到电流检测输入端，输出接负载；
    \item 电流计内阻LCR电桥测量，测量时自选方案消除接触电阻误差；
    \item 电流阶跃响应时间测量使用示波器采集输入和输出信号，以示波器显示为准。
\end{enumerate}

\subsection{充电控制部分}
\begin{enumerate}
    \item 充电控制部分测试使用台式电源供电，电流输入到充电控制输入端口，输出端口接负载仪。负载仪设置恒电阻模式，调节负载仪保证输出电压、电流在规定范围内；
    \item 对于实现恒流输出的充电控制部分，涉及到大功率DC-DC变换器，测试充电控制部分时电流从小到大设置，测试时控制电路耗散功率（输入功率-输出功率）不大于10W，如果确实可以大于10W操作，电路任何部分最高温度不得超过95摄氏度。在任何情况下，一旦超过耗散功率或温度限制立刻停止测试；
    \item 对于实现过流断电的充电控制部分，电流超调和响应时间满足要求的同时，应当完善保护电路工作逻辑——台式电源无法在断电之后及时停止输出。
\end{enumerate}

\section{评分标准}
\begin{table}[!htb]
    \centering
    \caption{考核评分标准（总分200分）}
    \resizebox{\linewidth}{!}{% 自适应表格宽度
    \begin{tabular}{llc}
        \toprule
        考核模块 & 考核要求 & 满分 \\
        \midrule
        \multirow{4}{*}{基础部分（50分）} & 电流计输入共模电压、总内阻符合基础要求 & 13 \\
        & 电流计量程、误差和响应时间符合基础要求 & 15 \\
        & 电流计使用分立元件搭建（运放等基础原件也可） & 5 \\
        & 控制电路具备过流保护功能且符合要求 & 15 \\
        \midrule
        \multirow{6}{*}{发挥部分（50分）} & 电流计内阻、误差满足发挥要求 & 10 \\
        & 充电控制电路具有恒流/恒压模式且输出电流达到2A & 15 \\
        & 转换效率达到要求（$>$90\%） & 2 \\
        & 转换效率优秀（$>$96\%）* & 12 \\
        & 具备短路保护等其他保护功能 & 5 \\
        & 设计美观（PCB设计与焊接）\&自由发挥 & 5 \\
        \midrule
        理论考核 & 对涉及知识点与相关问题进行现场提问 & 100 \\
        \midrule
        \multicolumn{2}{l}{合计} & 200 \\
        \bottomrule
    \end{tabular}
    }
    \note{*转换效率优秀：总转换效率大于96\%；使用集成芯片的，转换效率达到芯片标称值的90\%；若芯片未给出效率和输出电流的关系，效率不低于芯片标称最优值的90\%。}
\end{table}
